\documentclass[12pt,a4paper]{article}

% Codificação e idioma
\usepackage[utf8]{inputenc}
\usepackage[portuguese]{babel}

% Pacotes úteis
\usepackage{booktabs}
\usepackage{longtable}
\usepackage{graphicx}
\usepackage{hyperref}
\usepackage{geometry}
\usepackage{indentfirst}
\usepackage{float}

% Margens
\geometry{margin=2.5cm}


\newcommand{\projetotitulo}{Introdução à Manufatura Aditiva}
\newcommand{\grupoid}{Grupo Manufatura Aditiva}
\newcommand{\sprintnum}{3}
\newcommand{\sprintperiodo}{25/05 a 07/06/2026}
\newcommand{\scrummaster}{Lucas Moura}
\newcommand{\gitrepo}{https://github.com/viniciusnovelo/Introducao-a-Manufatura-Aditiva}
\newcommand{\gitcommit}{} % hash do commit desta entrega

\begin{document}

\begin{center}
{\Large\bfseries Relatório de Sprint \sprintnum}\\
{\large \projetotitulo}\\
\grupoid{} \ | \ \sprintperiodo{} \ | \ SM: \scrummaster
\end{center}

\section{Sprint Backlog}

\begin{longtable}{|c|p{3.8cm}|p{1.9cm}|p{2.2cm}|p{1.7cm}|p{4.2cm}|}
\toprule
ID & História & Estimativa & Responsável & Status & Critério de Aceitação \\
\midrule
1 & Criação do Instagram para Divulgação & Sprint 1 & Lucas Moura & Concluído & Perfil criado, configurado com identidade visual do projeto e contendo ao menos uma publicação inicial. \\ \hline

2 & Criação do formulário de divulgação & Sprint 1 & Vinícius & Concluído & Formulário criado, testado e apto para receber inscrições dos participantes. \\ \hline

3 & Início da divulgação para comunidade interna e externa & Sprint 1 & Vinícius & Concluído & Divulgação realizada em ao menos uma escola ou meio de comunicação definido pela equipe. \\ \hline

4 & Criação de slides para as aulas de SolidEdge para os calouros & Sprint 1 & Paulo & Concluído & Slides elaborados e revisados pela equipe antes da aplicação das aulas. \\ \hline

5 & Criação do quadro Kanban & Sprint 1 & Lucas Leão & Concluído & Quadro criado contendo colunas de acompanhamento e tarefas cadastradas. \\ \hline

6 & Coleta e organização de dados das inscrições dos alunos & Sprint 1 & Paulo & Concluído & Dados armazenados e organizados em planilha. \\ \hline

7 & Agendamento das salas e laboratórios que serão utilizados & Sprint 1 & Leonardo & Concluído & Reservas confirmadas junto à instituição para todas as aulas previstas. \\ \hline

8 & Verificação e manutenção das impressoras 3D e softwares no laboratório & Sprint 1 & Yuri & Concluído & Equipamentos testados e softwares funcionando corretamente antes do início do curso. \\ \hline

9 & Criação de relatório de Sprint 1 & Sprint 1 & Lucas Leão & Concluído & Relatório elaborado contendo atividades realizadas, resultados e pendências. \\ \hline

10 & Criação do repositório no Github & Sprint 1 & Vinícius & Concluído & Repositório criado com acesso para a equipe e estrutura inicial organizada. \\ \hline

11 & Treinamento da equipe nas ferramentas de modelagem 3D & Sprint 1 e 2 & Lucas Moura & Concluído & Equipe capacitada para executar operações básicas de modelagem no software. \\ \hline

12 & Elaboração dos materiais para as aulas 1 e 2 do curso & Sprint 2 & Equipe & Concluído & Materiais preparados e revisados antes da realização das aulas. \\ \hline

13 & Realização da aula 1 & Sprint 2 & Equipe & Concluído & Aula ministrada conforme cronograma e com registro de presença dos alunos. \\ \hline

14 & Realização da aula 2 & Sprint 2 & Equipe & Concluído & Aula ministrada conforme cronograma e com registro de presença dos alunos. \\ \hline

15 & Criação de relatório de Sprint 2 & Sprint 2 & Leonardo, Paulo, Yuri & Concluído & Relatório elaborado contendo atividades realizadas, resultados e pendências. \\ \hline

16 & Atualização do quadro Kanban & Sprint 2 & Lucas Leão & Concluído & Todas as tarefas da sprint atualizadas conforme o andamento das atividades. \\ \hline

17 & Atualização do repositório Github & Sprint 2 & Vinícius & Concluído & Materiais e documentos da sprint enviados e organizados no repositório. \\ \hline

18 & Treinamento da equipe nas ferramentas de impressão 3D & Sprint 3 & Lucas Moura & Não iniciado & Equipe apta a operar as impressoras e configurar impressões básicas. \\ \hline

19 & Elaboração dos materiais para as aulas 3 e 4 do curso & Sprint 3 & Equipe & Concluído & Materiais preparados e aprovados pela equipe antes das aulas. \\ \hline

20 & Realização da aula 3 & Sprint 3 & Equipe & Concluído & Aula ministrada conforme cronograma e com registro de presença dos alunos. \\ \hline

21 & Criação de relatório de Sprint 3 & Sprint 3 & Leonardo, Paulo, Yuri & Concluído & Relatório elaborado contendo atividades realizadas, resultados e pendências. \\ \hline

22 & Atualização do quadro Kanban & Sprint 3 & Lucas Leão & Não iniciado & Todas as tarefas da sprint atualizadas conforme o andamento das atividades. \\ \hline

23 & Atualização do repositório Github & Sprint 3 & Vinícius & Em andamento & Materiais e documentos da sprint enviados e organizados no repositório. \\ \hline

24 & Elaboração dos materiais para as aulas 5 e 6 do curso & Sprint 4 & Lucas Moura & Não iniciado & Materiais finalizados, revisados e disponíveis para aplicação das aulas. \\ \hline

25 & Realização da aula 4 & Sprint 4 & Equipe & Não iniciado & Aula ministrada conforme cronograma e com registro de presença dos alunos. \\ \hline

26 & Realização da aula 5 & Sprint 4 & Equipe & Não iniciado & Aula ministrada conforme cronograma e com registro de presença dos alunos. \\ \hline

27 & Criação de relatório de Sprint 4 & Sprint 4 & Leonardo, Paulo, Yuri & Não iniciado & Relatório elaborado contendo atividades realizadas, resultados e pendências. \\ \hline

28 & Atualização do quadro Kanban & Sprint 4 & Lucas Leão & Não iniciado & Todas as tarefas da sprint atualizadas conforme o andamento das atividades. \\ \hline

29 & Atualização do repositório Github & Sprint 4 & Vinícius & Não iniciado & Materiais e documentos da sprint enviados e organizados no repositório. \\ \hline

30 & Realização da aula 6 & Sprint 5 & Equipe & Não iniciado & Aula final realizada e conteúdo previsto concluído. \\ \hline

31 & Criação de formulário para coleta de feedback & Sprint 5 & Leonardo, Yuri & Não iniciado & Formulário criado contendo perguntas objetivas e discursivas sobre o curso. \\ \hline

32 & Coleta de feedback dos alunos através de formulário & Sprint 5 & Leonardo, Yuri & Não iniciado & Respostas coletadas de parte significativa dos participantes do curso. \\ \hline

33 & Criação de relatório de Sprint 5 & Sprint 5 & Leonardo, Paulo, Yuri & Não iniciado & Relatório elaborado contendo atividades realizadas, resultados e pendências. \\ \hline

34 & Atualização do quadro Kanban & Sprint 5 & Lucas Leão & Não iniciado & Todas as tarefas da sprint atualizadas conforme o andamento das atividades. \\ \hline

35 & Atualização do repositório Github & Sprint 5 & Vinícius & Não iniciado & Materiais e documentos da sprint enviados e organizados no repositório. \\ \hline

36 & Criação do relatório final & Sprint 5 & A definir & Não iniciado & Relatório final concluído contendo descrição completa do projeto, resultados e conclusões. \\

\bottomrule
\end{longtable}

\section{Burndown Chart}
\begin{figure}[h!]
\centering
\includegraphics[width=1.0\textwidth]{Imagens/burndown_sprint\sprintnum.png}}
\caption{Burndown -- Sprint \sprintnum}
\end{figure}

% Análise do burndown (mínimo 3 linhas):
O gráfico acima representa o engajamento e o trabalho efetivo realizado pelo grupo durante a Sprint 3, sendo que a linha azul representa o progresso ideal no decorrer de cada dia e a linha laranja o progresso realizado efetivamente. Para cada atividade da Sprint 3 foi atribuído um valor de esforço representado no eixo vertical do gráfico, além disso, nessa Sprint duas atividades planejadas não foram concluídas.

\section{Artefatos Produzidos}

Em continuidade aos materiais desenvolvidos nas semanas anteriores, realizou-se a confecção de slides para o uso durante as aulas 3 e 4. Os mesmos continham explicações breves e elementos visando facilitar a explicação do instrutor e o entendimento dos alunos em relação aos conhecimentos ensinados. Também foram elaboradas peças detalhadas para prática dos alunos no Solid Edge.
\begin{figure}[H]
\centering

\begin{minipage}{0.48\textwidth}
    \centering
    \includegraphics[width=\textwidth]{Imagens/Slide - Montagem.png}
    \caption{Slide - Montagem}
\end{minipage}
\hfill
\begin{minipage}{0.48\textwidth}
    \centering
    \includegraphics[width=\textwidth]{Imagens/Peça.png}
    \caption{Peça - Prática de modelagem}
\end{minipage}

\vspace{2.0cm}

\begin{minipage}{0.4\textwidth}
    \centering
    \includegraphics[width=\textwidth]{Imagens/Aula 30_05 - 1.jpeg}
        \caption{Aula 30/05}
\end{minipage}
\hfill
\begin{minipage}{0.4\textwidth}
    \centering
    \includegraphics[width=\textwidth]{Imagens/Aula 30_05 - 2.jpeg}
        \caption{Aula 30/05}
\end{minipage}
\end{figure}

\section{Atualização com o Parceiro}
O contato com os parceiros ocorreram no decorrer das semanas no grupo do WhatsApp e na aula do dia 30/05. Nessas interações foram repassadas orientações, avisos, e propriamente o conteúdo nos períodos de aulas. Vale ressaltar que a quarta aula do curso foi postergada para o dia 13/06.

\section{Participação do Calouro}
Durante o encontro presencial do dia 30/05, dois dos calouros, Isaac e Otávio, atuaram como instrutores ministraram o decorrer da aula. Paralelamente aos encontros presenciais em si, os calouros auxiliaram no desenvolvimento dos materiais e organização dos conteúdos para os alunos.

\section{Impedimentos e Desvios}
O grande impedimento em relação a essa Sprint foi a expiração repentina das licenças do Solid Edge que impossibilitou os alunos de acessarem o software no dia 30/05, entretanto, o grupo decidiu prosseguir com a aula e os conteúdos programados.

Ademais a isto, problemas pessoais e de saúde impossibilitaram o SM, que seria responsável pelo treinamento nas ferramentas de impressão 3D, de realizar o treinamento. 

\section{Retrospectiva}
\subsection*{O que funcionou bem}
Mesmo com os impedimentos apresentados do dia 30/05, os alunos se mostraram engajados e sem muitas dificuldades ou dúvidas em relação aos conteúdos lessionados.

\subsection*{O que precisa melhorar}
A organização das próximas etapas precisa ser precisa e bem executada para evitar atrasos e entregas vagas.

\subsection*{Ação concreta para a próxima sprint}
Realizar o treinamento da equipe na impressão 3D o mais breve possível e garantir que todas as ferramentas pertinentes ao andamento do curso estejam funcionando de forma adequada.

\section{Planejamento da Próxima Sprint}
Na Sprint 4 a equipe dará prosseguimento com as aulas 4 e 5 do minicurso e com a preparação dos materiais adequados para ministrar as mesmas. Além disso, os membros do grupo tem como objetivo primordial verificar os processos operacionais para manejo das impressoras 3D e posterior capacitação dos calouros.

Nesta Sprint, será feito o contato com a professora Michele para garantir a emissão dos certificados para os alunos, que deverão ser disponibilizados na sexta aula.
\end{document}